\section{Gauge-Theoretic Encoding of Quantum States}
\label{sec:gauge-construction}

\subsection{Definition of the Model}

We consider a one- or two-dimensional lattice with a $\mathbb{Z}_8$ gauge theory. The fundamental variables are $\mathbb{Z}_8$ fluxes $f_e \in \{0,1,\dots,7\}$ assigned to each edge $e$ of the lattice. We use additive notation for the group operation.

At each vertex $v$, we impose the Gauss law constraint:
\begin{equation}
\sum_{e \ni v} \eta_{v,e} f_e \equiv \rho_v \pmod{8},
\end{equation}
where $\rho_v \in \mathbb{Z}_8$ is a \textit{defect charge} at vertex $v$ and $\eta_{v,e} \in \{-1,+1\}$ encodes the orientation of edge $e$ with respect to the vertex. When $\rho_v = 0$ for all $v$, we recover the standard gauge-invariant sector. A non-zero $\rho_v$ represents a topological defect.

The gauge group $\mathbb{Z}_8$ acts by shifting fluxes along cycles: for any closed loop $\gamma$, we can shift all edges in $\gamma$ by a constant $g \in \mathbb{Z}_8$. Gauge-invariant observables are Wilson loops:
\begin{equation}
W_\gamma = \exp\left(\frac{2\pi i}{8} \sum_{e \in \gamma} f_e\right).
\end{equation}

\subsection{Dictionary: Qubits to Fluxes}

Qubit degrees of freedom are encoded in oriented flux differences across each vertex. In a 1D chain, vertex $v_i$ has a left incident edge $e_{i-1}$ and a right incident edge $e_i$, so the local Gauss law becomes
\begin{equation}
f_{e_i} - f_{e_{i-1}} \equiv \rho_i \pmod{8}.
\end{equation}
For open boundary conditions we set $f_{e_0} = f_{e_n} = 0$.

The logical Pauli operators may be represented on the slice flux basis $\{|q\rangle : q \in \mathbb{Z}_8\}$ by:
\begin{align}
Z_i |q\rangle &= (-1)^q |q\rangle, \\
X_i |q\rangle &= |q + 4 \pmod{8}\rangle, \\
Y_i &= i X_i Z_i.
\end{align}

In this representation, Clifford gates act by conjugating these Pauli operators and therefore preserve the Gauss law, while $T$ gates change the defect charges $\rho_i$.

\subsection{Clifford Gates as Gauge Transformations}

Clifford gates correspond to gauge transformations that preserve the Gauss law constraint:

\begin{itemize}
    \item \textbf{Hadamard (H)}: Exchanges X and Z bases, realized as a linear transformation of edge fluxes that preserves the Gauss law.
    \item \textbf{Phase (S)}: Multiplies by a phase depending on flux, does not create defects.
    \item \textbf{CNOT}: Transforms fluxes on two vertices while preserving the constraint.
\end{itemize}

\begin{lemma}
All Clifford operations map gauge-invariant states to gauge-invariant states. They act as unitary transformations within the physical Hilbert space.
\end{lemma}

\begin{proof}
Clifford gates preserve the stabilizer group, which corresponds to the Gauss law constraints. Thus they map valid flux configurations to valid flux configurations.
\end{proof}

\subsection{T-Gates as Topological Defects}

The $T$ gate is defined as $T = \text{diag}(1, e^{i\pi/4})$ in the computational basis. In our gauge theory, a $T$ gate on qubit $i$ corresponds to inserting a \textbf{topological defect} with charge $\tau = 1$ (mod 8) at vertex $v_i$.

Formally, the action of $T$ on a flux configuration is:
\begin{equation}
T_i | \{f_e\} \rangle = e^{i \frac{\pi}{4} q_v(\{f_e\})} | \{f_e\} \rangle
\end{equation}
where $q_v(\{f_e\}) = 1$ if the Gauss law is violated at $v_i$ (i.e., a defect is present), and $0$ otherwise.

A defect is characterized by:
\begin{itemize}
    \item Position: vertex $v_i$
    \item Charge: $\tau_i \in \mathbb{Z}_8$ (the amount of Gauss law violation)
\end{itemize}

The state after applying a sequence of $T$ gates is:
\begin{equation}
|\psi\rangle = \sum_{\{f_e\}} \exp\left(i \sum_{v \in D} \frac{\pi}{4} \tau_v\right) |\{f_e\}\rangle
\end{equation}
where $D$ is the set of defect positions.

\subsection{State as Defect Configuration}

After a Clifford+T circuit, the quantum state is completely specified by:
\begin{enumerate}
    \item The lattice geometry (fixed once and for all)
    \item The defect configuration $D = \{(v_i, \tau_i)\}$ is specified by positions and charges
\end{enumerate}

The description size is $O(|D|)$, where $|D|$ is the number of $T$ gates applied. For a circuit of depth $d$ with $O(1)$ $T$ gates per layer, $|D| = O(nd)$.
